%!TEX root = TheoretischePhysikC.tex

\section{Vorlesung: Hilbertraum}
Um die Zustände der Quantenmechanik zu beschreiben wird der \emph{Hilbertraum} eine größere Rolle spielen. Um diesen mathematischen Raum zu definieren, benötigen wir erst \emph{komplexe Vektorräume}. Solche Räume werden analog mit dem reellen Fall definiert. 

Ein \emph{komplexer Vektor} $\mbf{v}\in \mathbb{C}^n$ ist eine Spalte von $n$ komplexen Zahlen
\begin{equation}
	\mbf{v} = 
	\begin{pmatrix}
		v_1 \\ v_2 \\ \vdots \\ v_n
	\end{pmatrix}.
\end{equation}
Es gelten die folgenden Rechenoperationen
\begin{enumerate}
	\item Multiplikation mit Skalar $a\in\mathbb{C}$
	\begin{equation}
		a\mbf{v} = 
	\begin{pmatrix}
		av_1 \\ av_2 \\ \vdots \\ av_n
	\end{pmatrix}.
	\end{equation}
	\item Addition
	\begin{equation}
		\mbf{v}+\mbf{w} = \begin{pmatrix}
		v_1 \\ v_2 \\ \vdots \\ v_n
	\end{pmatrix} + \begin{pmatrix}
		w_1 \\ w_2 \\ \vdots \\ w_n
	\end{pmatrix} =
	\begin{pmatrix}
		v_1+w_1 \\ v_2+w_2 \\ \vdots \\ v_n+w_n
	\end{pmatrix}.
	\end{equation}
	\item Kreuz-Operation
	\begin{equation}
		v^\dag = \begin{pmatrix}
		v_1 \\ v_2 \\ \vdots \\ v_n
	\end{pmatrix}^\dag \equiv \begin{pmatrix}
		v_1^* & v_2^* & \cdots & v_n^*
	\end{pmatrix}
	\end{equation}
	\item Skalarprodukt
	\begin{equation}
		\mbf{v}^\dag\mbf{w} = \begin{pmatrix}
		v_1^* & v_2^* & \cdots & v_n^*
	\end{pmatrix} \begin{pmatrix}
		w_1 \\ w_2 \\ \vdots \\ w_n
	\end{pmatrix} =
	\sum_{j=1}^n v_j^*w_j.
	\end{equation}
\end{enumerate}
Der Raum $V$ aller solcher Vektoren mit festem $n$ wird als \emph{komplexer Vektorraum} bezeichnet. 

Hierbei wird $n$ als \emph{Dimension} von $V$ definiert. Oft verwendet wird die Notation $\mathbb{C}^n$ für den oben definierten Vektorraum $V$.

\begin{beispiel}
	Der Vektorraum $\mathbb{C}^2$ besteht aus allen zweidimensionalen Vektoren
	\begin{equation}
		\mbf{v} = 
		\begin{pmatrix}
			v_1 \\ v_2
		\end{pmatrix}.
	\end{equation}
	Seien
	\begin{equation}
		\mbf{v} = \begin{pmatrix}
			1+i \\ 1-i
		\end{pmatrix} \quad \text{und} \quad \mbf{w} = \begin{pmatrix}
			1 \\ -1
		\end{pmatrix}
	\end{equation}
	zwei Vektoren. Deren Skalarprodukt lautet
	\begin{equation}
		\mbf{v}^\dag \mbf{w} = \begin{pmatrix}
			1-i & 1+i
		\end{pmatrix} \begin{pmatrix}
			1 \\ -1
		\end{pmatrix} = -2i.
	\end{equation}
\end{beispiel}

Unendlich dimensionale Vektorräume sind ebenfalls möglich und spielen eine sehr wichtige Rolle in der Beschreibung quantenmechanischer Systeme. Wie sind solche Vektorräume zu verstehen? Man kann sich unendlich dimensionale Vektoren als Spalten mit unendlich vielen Einträge vorstellen, z.B.:
\begin{equation}
	\mbf{v} = \begin{pmatrix}
		v_1 \\ v_2 \\ \vdots \\
	\end{pmatrix}
\end{equation}
Alle Rechenregeln für Vektorräume gelten auch in den unendlichen dimensionalen Fällen. Das Skalarprodukt $\mbf{v}^\dag \mbf{w}$ wird wie folgt berechnet
\begin{equation}
	\mbf{v}^\dag\mbf{w} =
	\sum_{j=1}^\infty v_j^*w_j.
\end{equation}
Das heißt, in der Regel, dass man in allen Gleichungen $n=\infty$ setzt. Es gibt mehrere subtile Aspekte, die man in Betracht ziehen müsste, um rigorös mit unendlich dimensionalen Vektorräumen zu arbeiten. Allerdings werden wir in dieser Vorlesung diese Aspekte nur nebenbei erwähnen.

Für uns ist das wichtigste Beispiel einen unendlich dimensionalen Vektorraumes der Raum der \emph{quadratintegrierbaren Funktionen} (der als $L^{2}(\mathbb{R})$ bezeichnet wird) welcher $\mathbb{R}$ auf $\mathbb{C}$ abbildet. Hierbei bedeutet quadratintegrierbar, dass
\begin{equation}
	\int_{-\infty}^\infty |f(x)|^2 \, dx < \infty.
\end{equation}
Quadratintegrierbare Funktionen können addiert, subtrahiert, und mit Skalargrößen multipliziert werden. Bezüglich dieser Operationen bleiben die Ergebnisse immer quadratintegrierbar. Daher bildet die Menge aller quadratintegrierbaren Funktionen einen abstrakten Vektorraum. Man betrachtet eine quadratintegrierbare Funktion $f(x)$ als einen Vektor. 

\begin{beispiel}
	Ein \emph{Nicht}-Beispiel von einer quadratintegrierbaren Funktion wird von der konstanten Funktion $f(x) = 1$ geliefert:
	\begin{equation}
		\int_{-\infty}^\infty |f(x)|^2\, dx = \int_{-\infty}^\infty 1\, dx = \infty.
	\end{equation}
\end{beispiel}


